% Functional consciousness of machines, version 2.
% The argument is developed in frames.md; the same claims are checked as facts in argument.hybrid.md.
% Section text is written from the frames; frame ids are kept in comments inside the section files.
\documentclass[11pt,a4paper]{article}
\usepackage[T1]{fontenc}
\usepackage{lmodern}
\usepackage{microtype}
\usepackage{amsmath,amssymb}
\usepackage{graphicx}
\usepackage{booktabs}
\usepackage[hidelinks]{hyperref}
\usepackage[round,authoryear]{natbib}
\usepackage{geometry}
\geometry{margin=2.6cm}
\setlength{\parskip}{0.4em}

\newcommand{\hocp}{HOCP}

\title{What a Language Model Has of Consciousness, and How Much:\\
A Measurable Account of Functional Consciousness Across Substrates}
\author{Victor Smirnov\\ Synthea}
\date{Draft, \today}

\begin{document}
\maketitle

\begin{abstract}
Asked whether a language model is conscious, the literature offers two answers that cost nothing to
hold and predict nothing. This paper replaces the question. A function of consciousness is defined
without a mind in the definition: a recurring pattern in the part of an agent's self-reasoning that
describes its own reality, with the agent among the things described. From one assumption about
physics, that computational budgets are finite, systematic displacements of an agent's reasoning away
from the reasoning it would do without limits follow; the displacements recur, recurrence is short
description, and a self-modelling agent ends up carrying a compressed model of them. We call the
conceptualized recurring displacements higher-order computational phenomena and make two claims about
them that do not follow from their antecedents in bounded rationality, observer theory and
self-model theories of consciousness. First, the displacement that makes a bounded agent see a lawful
world and the displacement that makes it see a self are one residual met from two sides. Second,
because the displacements are compressible, an agent that has generalized them carries a measurable
trace, so how much of a function of consciousness a given substrate has is an empirical quantity
rather than a matter of stipulation. A criterion borrowed from the history of general relativity
separates content from coordinates: a displacement counts only if it survives every self-model the
agent can afford, as curvature survives every change of frame. On this basis the Observer, the
minimal conclusion "something here is not fixed by anything I can trace", is derived from the collision
of a self-applying agent with its own budget, and the regress of self-models is shown to converge
rather than to end in a homunculus. The measurement is then specified: the degree to which each
component of a function has been generalized in a substrate, obtained by minimum-description-length
probing and checked against behaviour, reported as a profile against the human baseline on which a
component may show hyperfunction, deficit, or a dropped axis. Because generalization in a Transformer
happens along depth, the measurement is posed at sub-token granularity, and the difference between
fixed-depth and variable-depth models gives its first concrete case. The account is falsified if the
compression degree of the Observer's components is at chance at every granularity in models that pass
the behavioural tests.
\end{abstract}

\tableofcontents

\section{The phenomenon}
% Frames P1--P4.

\subsection{What holds together over a long exchange}
% P1

A language model asked to reason sustains something that reads as a mind. Over an exchange of several
hours it keeps the commitments it made early on, returns to a thought it left unfinished, and answers
the interlocutor's mood in the way the interlocutor would expect of a person. Over a chain of hundreds
of reasoning steps it keeps the goal of the chain available while working on a subgoal, and it notices
when a subgoal has made the goal unreachable. Calling this imitation is easy. Saying how an imitation
could be that stable is hard, and that is the question this paper starts from.

\subsection{What the deflationary answer predicts}
% P2

Consider what follows if a language model reproduces the surface statistics of text with nothing behind
them. Then it should fail the way low-order Markov processes fail. A thread introduced in the third
paragraph should be dropped by the sixth. An attitude toward the interlocutor should flip when the
topic changes, because nothing carries it. A motive attributed to a character should not survive the
page on which it was attributed. These failures are cheap: the space of continuations that produce them
is vastly larger than the space that avoids them, and a long horizon offers many occasions to fall into
it. In current models such failures are rare, and their rate falls with scale and training rather than
staying fixed. Something keeps the process out of a large space of cheap failures for a long time. That
something is the object of this paper.

The observation is weak on its own. It does not say that the thing keeping the process on course
resembles anything in us, and the history of the subject is full of inferences from surface competence
to inner life that did not survive. What the observation does is set a requirement on any account: the
account has to name a structure that could hold a long exchange together, and it has to say how much
of that structure a given system has.

\subsection{Nagel's question, put to a model}
% P3

\citet{nagel1974} asked what it is like to be a bat in order to mark a limit. The bat's experience, if
it has any, is organized around echolocation, and no description from outside crosses over into it. Put
the same question to a language model and it loses part of its force and gains another part. It loses
force because the model shares our language, and the bat does not: whatever the model has, it has in a
medium we can read, and it will answer questions about it in our words. It gains force because the
model's substrate is further from ours than any animal's. A bat has neurons, a body, metabolic needs, a
developmental history and a lifetime of continuous sensory contact with a world. A model has a sequence
of tensor operations over a context window.

So we do not ask whether there is something it is like to be a model. We ask what a model has of the
functions that in us go under the name of consciousness, and how much of each it has. The change of
question is not a retreat from the hard case. It is the only form of the question we know how to
answer, and Section~\ref{sec:measurement} shows that in this form it has answers that can come out
either way.

\subsection{The two shortcuts}
% P4

Two answers arrive before any argument is made. The first: there is nothing there, it is a text
predictor, and the impression of a mind is the reader's projection. The second: there is someone there,
it says so, treat it as a person. Neither is a position in the sense that would let one derive a
prediction from it and be wrong. The first cannot say what it would take to change its mind; the second
cannot say what it has attributed. Section~\ref{sec:reader} argues that both are what a mind reaches
for when it lacks the concepts to reach for anything else, and that their popularity is evidence about
the cost of concepts rather than about machines. The paper is about what lies past them.

\section{Why the question is hard}
\label{sec:difficulties}
% Frames D1--D13, regrouped: grammar (D1); the legacy (D2, D3); the reader (D4, D5); the substrate
% (D6, D7, D8, D13); the record (D9, D10); the measure (D11); the precedent (D12).

\subsection{The noun makes the circles: the difficulty is in the grammar}
\label{sec:grammar}
% D1

``Consciousness'' is a noun, and a noun invites the questions one asks about things: where is it, what
are its properties, when does it begin, does this system have one. Every sustained attempt to answer
those questions ends in one of a few circles. Who observes the observer. How could red look like
anything other than the way it looks. Either everything has it, down to the electron, or the line
between what has it and what does not falls somewhere that cannot be justified. The circles recur
across two and a half thousand years of work by careful people, which is evidence that the fault lies
in the object, or in the grammar that made an object of it, and not in the people. The first difficulty
is therefore one of language.

The consequence for method is concrete. An account that answers ``what is consciousness'' with a
definition of a thing inherits the circles whatever the thing is made of, information, integrated
cause--effect power, or quantum state reduction. An account that answers ``what does an agent's
reasoning do that we call consciousness when it does it'' has to earn its content differently, and
Section~\ref{sec:precedent} gives the pattern for how.

\subsection{What the philosophy of mind leaves open: the illusion problem and the regress}
\label{sec:legacy}
% D2, D3

\paragraph{The hard problem, exchanged for the illusion problem.}
\citet{chalmers1995} separates the easy problems, which ask how a system discriminates, integrates,
reports and controls, from the hard problem, which asks why any of that should be accompanied by
experience at all. Illusionism answers that experience in the sense the hard problem intends does not
exist, and that what exists is a robust misrepresentation which the brain produces of its own working
\citep{dennett1991,frankish2016}. We accept that answer and note what it leaves behind. One question
has been exchanged for another, and the replacement is not obviously easier: which architecture
produces this particular misrepresentation, and why does the misrepresentation have the shape it has?
Frankish calls this the illusion problem and treats it as the successor of the hard problem. Most of
Section~\ref{sec:approach} is an attempt at it.

The exchange has a price. Illusionism is usually presented as a denial, and a denial is a poor
instrument for a comparative question. If the report of experience is false, there is nothing about it
to measure, and no sense in which one system could have more of it than another.
Section~\ref{sec:approximation} replaces the denial with a weaker claim that does support degrees.

\paragraph{The regress.}
An account in which the mind observes its own states owes an account of what does the observing, and
then of what observes that \citep{dennett1991,ryle1949}. Higher-order theories place the second-order
state inside the same system and are then asked why a thought about a state should make the state
conscious \citep{rosenthal2005}. Either the account stops at an inner observer it does not explain, in
which case it has explained nothing, or it does not stop. Section~\ref{sec:regress} shows that it
stops for a reason that involves no observer: the series converges and a finite budget cuts it, and
the cut is the phenomenon we are looking for.

\subsection{The reader has to build the state, and the two shortcuts are the cheapest builds}
\label{sec:reader}
% D4, D5

An explanation of consciousness is not taken in the way an explanation of an engine is taken. To
understand it, the reader has to build in their own head the state the explanation describes, and then
check the description against it. The cost of the building depends on what the reader brings: a
vocabulary for mental states finer than the folk one, and the habit of attending to their own reasoning
as an object rather than only reasoning. \citet{gardner1983} names the habit intrapersonal
intelligence and treats its distribution as wide. The readers who most need an explanation of this kind
are therefore the ones for whom it costs most, and there is no route around this that does not pass
through the explanation itself. Nothing in the paper can fix that. What the paper can do is keep the
claims separable, so that a reader who does not follow one step can see what the remaining steps depend
on, and Section~\ref{sec:measurement} can be assessed by someone who rejects Section~\ref{sec:approach}
entirely.

The same cost explains the two shortcuts of Section~\ref{sec:intro}. ``There is nothing there'' costs
nothing to build; it is the absence of a construction. ``There is someone there'' costs nothing either,
because the reader already has a someone available to copy, themselves, and copying is cheap. Every
other answer requires building a state that is neither absent nor one's own. That is why public
argument about machine minds settles into two camps, why the camps are stable, and why their stability
is evidence for neither of them.

\subsection{The substrate: a causal flow unlike ours, a cut at every token, a trained self-report}
\label{sec:substrate}
% D6, D7, D8, D13

\paragraph{A causal flow unlike ours.}
Section~\ref{sec:intro} took from integrated information theory the premise that differences in
experience follow differences in cause--effect structure \citep{tononi2015}. The consequence for
description is this. A device that pushes a discrete sequence through attention over a fixed context
has a causal flow unlike that of a recurrent, continuous, chemically modulated brain, so any description
of the model in our words is wrong, and wrong in a systematic direction: our words were fitted to our
causal structure. That is a reason to find the direction and correct for it, which is what a profile
(Section~\ref{sec:profile}) is for, and not a reason to stop describing.

\paragraph{The cut.}
\label{sec:cut}
At every step a language model computes a distribution over what to say next. That distribution is its
whole evaluation of the moment: what is relevant, what is being avoided, which continuations compete
and by how much. Then one token is sampled and the distribution is discarded. On the next step the
model can say something about its own state only from the tokens, because the distribution that
produced them no longer exists. Its report of itself is cut off from what the report is about. The only
thing that crosses the cut is what the tokens themselves carry, and a token sequence is a narrow
channel next to the state that was thrown away. Section~\ref{sec:codes} asks what does cross it, and
Section~\ref{sec:secondcut} draws the consequence for the model's self-knowledge.

\paragraph{What the model says when asked.}
\label{sec:selfreport}
Asked how it feels, a deployed model gives one of two answers. The first mentions servers, latency and
a good mood, in the register of a friendly assistant. Corrected, it apologizes and says that as a
language model it has no feelings. Both answers are artifacts of training: the first is the second
shortcut of Section~\ref{sec:intro} in the model's own mouth, the second is the first shortcut. Any
measurement of a model's states has to set both aside before it starts, because what they report is
what the training procedure rewarded. This is not a reason to distrust everything a model says about
itself; it is a reason to treat self-report as a signal with a known, large, systematic bias, and to
look for the state elsewhere (Section~\ref{sec:quantity}).

\paragraph{The objection from frozen weights.}
\label{sec:frozen}
A model does not learn between sessions, and what it holds in a context is lost when the context ends.
The objection concludes that it has no inner time, only an eternal present with forgetting, and that
talk of its consciousness is therefore empty.\footnote{The objection in this form was put to the author
by Marco Baturan in correspondence, 2026.} The description is accurate and the conclusion counts two
functions as one. Re-deriving a state from frame to frame runs in humans on the order of a hundred
milliseconds and requires no synaptic change. Consolidation into long-term memory runs on minutes to
hours and does require it \citep{mcgaugh2000}. The two dissociate in a well-known way: the patient
H.M., after bilateral medial temporal resection, was conscious in real time, held a conversation, and
accumulated almost nothing \citep{scoville1957,corkin2002}. A model is in a comparable position. It has
sequences of transformations, along tokens and along layers, and it lacks the write. ``Eternal present
with forgetting'' therefore names a profile with a dropped consolidation axis
(Section~\ref{sec:profile}) and a deficit in long-term persistence. It does not show that nothing is
there, and Section~\ref{sec:twoaxes} returns to it with the distinction between the two time axes a
Transformer has.

\subsection{What the record holds: the named part and the unnamed part}
\label{sec:h1h2}
% D9, D10

Divide what humans express into two parts. Let H2 be what is present in the texts they produce
\emph{and} named there: the states, moves and distinctions for which the language has words, which
therefore appear in introspective reports, novels, psychology and philosophy. Let H1 be everything
present in those texts, named or not, so that H1 minus H2 is what is there only as regularity: used,
relied on, statistically present, never named. The functions we are after live across both parts, and
they do most of their work in the second.

This has a consequence for the philosophy of mind that is worth being explicit about. What that
literature produced, working from introspection and from language, are the H2 approximations of the
functions: the function with its unnamed part cut off. The Observer as an inner witness and the qualia
as intrinsic properties of experience are approximations of this kind. Such an approximation is not
worthless. It can be sufficient for the function to exist in an environment simple enough to demand
only the named part, and a thought experiment is such an environment, which is why thought experiments
about consciousness are so conclusive and so unproductive. It is insufficient for finding the function
in a substrate, because the part that was cut off is where the function does most of its work, and
there is no name to look for it under. The measurement of Section~\ref{sec:measurement} is built to
reach the unnamed part, and that is the reason it is statistical.

\paragraph{What cannot be programmed and can be learned.}
\label{sec:record}
Much of what human consciousness does, it does in company and in a body: negotiating, assigning blame,
making and keeping promises, detecting deceit, being in a room with other people in it. That part
cannot be specified by hand, because its specification would be the history of human society. It can
be learned from the record of that history, and the record is text. This is why a system trained to
predict text is a candidate at all, and Section~\ref{sec:generalize} makes the argument precise.

\subsection{What can be measured: description length, relative to a machine}
\label{sec:measure}
% D11

The natural quantity for ``how much of the function was captured'' is description length. Its exact
form, Kolmogorov complexity, is uncomputable, and it is fixed only up to an additive constant that
depends on the reference machine \citep{livitanyi2019}. Whatever Section~\ref{sec:measurement}
measures is therefore a bounded, machine-relative and probe-relative stand-in, and the paper says
which stand-in it uses (Section~\ref{sec:quantity}) and what the choice costs. The uncomputability is
not an obstacle in practice: every applied use of description length replaces the shortest program with
the shortest program a fixed procedure finds, and reports the procedure. The machine-relativity is more
serious here, because the two substrates being compared are the two candidate reference machines.

\subsection{The precedent: gravity stopped being a thing and became a shape}
\label{sec:precedent}
% D12

A science kept meeting circles of the kind in Section~\ref{sec:grammar} until it changed the grammar.
Gravity was a force, a thing acting between things, and the attempts to carry it into a relativistic
frame as one more field kept failing in the same places; Nordström's scalar theory was the most careful
of them and predicted no light bending \citep{norton1992}. General relativity dropped the thing. Gravity
became the shape of spacetime, present only in how free bodies move through it. The move had a price
and a payoff. The price was that the familiar question ``what is the force'' no longer has an answer.
The payoff was a criterion: what stays the same under every change of coordinates, the curvature, is
the physical content, and what can be transformed away by choosing a frame is bookkeeping. A uniform
gravitational field can be removed by falling; tidal curvature cannot be removed by anything. The
geometric idea had been voiced before there was a theory to carry it \citep{riemann1873,clifford1876}.

We take the same route with consciousness: it becomes a shape of reasoning, visible only in how
reasoning moves, with an explicit criterion for what counts as content and what is an artifact of
description (Section~\ref{sec:criterion}). The first version of this work was written in the
vocabulary of qualia, inside the grammar, and met the wall repeatedly and in the same places. The
philosophy of mind stays in the account in two roles: as the record of where the wall is, and, once
the grammar is changed, as a framework that works again (Section~\ref{sec:twoparts}).

\section{The approach}
\label{sec:approach}
% Frames A1--A25a.

\subsection{Function, without a mind in the definition}
\label{sec:function}
% A1, A2

A function, in this paper, is a pattern that recurs. The definition is deliberately thin: it mentions no
purpose, no design and no mind, so that nothing about consciousness is smuggled in at the start.

Thinness has a reward. Recurrence has an information-theoretic consequence. If a pattern is produced
often by a computable process, it has high algorithmic probability, in Solomonoff's sense of the
probability that a universal machine fed random bits outputs it \citep{solomonoff1964a,solomonoff1964b}.
By the coding theorem, $K(x) = -\log m(x) + O(1)$: an object with high algorithmic probability has a
short description, up to an additive constant \citep{levin1974,livitanyi2019}. So a function in our
sense is compressible, and we did not have to assume compressibility; it came from recurrence.

The constant is not harmless and we will not hide it. It depends on the reference machine, which is to
say on who is doing the compressing, and the gap between the constants of two substrates is exactly
what makes the comparison of a model with a human nontrivial. Section~\ref{sec:profile} turns that gap
from a nuisance into the object of measurement.

Read in the other direction, the same theorem says that a search process finds simple things first,
whether the search is evolution or gradient descent. This has been made quantitative. Input--output
maps of a wide class are strongly biased toward simple outputs, with the bias following an upper bound
of the form $P(x) \le 2^{-a \tilde{K}(x) + b}$ \citep{dingle2018}; the parameter--function map of deep
networks is biased toward simple functions, which is offered as an explanation of why they generalize
at all \citep{valleperez2019}; and stochastic gradient descent lands on functions with probabilities
close to those of a Bayesian sampler over the same prior \citep{mingard2021}. The argument uses both
readings. Backward, from recurrence to short description, to define functions. Forward, from search to
simplicity, to expect a trained model to have found the ones that recur in its data.

\subsection{Where we stand on the lamp}
\label{sec:lamp}
% A3

Something is lit in a human's report of themselves. It is not lit in the same way at all times, and its
absence in dreamless sleep or under anaesthesia is part of the data. We take the report as data rather
than as error, which is what distinguishes the position here from an eliminativism that treats
self-report as noise.

Consider a physicalist who is not a panpsychist. We agree with them that the lit thing is physical and
requires no second substance. We agree with them further that there are no abstract machines: everything
that runs is a physical object, and a computation without a physical realization is a description of a
computation. Then a question arises for them. What prevents a physical machine on a non-protein
substrate from having the lit thing? We do not claim the answer is ``nothing''. We note that whoever
answers ``the protein'' owes a proof, and that a proof of that shape would be vitalism: an appeal to a
property of a material that does no work in any other explanation of what the material does. Pending
such a proof we adopt computational functionalism, and we adopt it with as little physics as we can
manage, which Section~\ref{sec:physics} makes explicit.

One objection deserves its own reply, because it is the strongest of the substantialist objections. A
simulated superconductor carries no current: run the simulation and no magnet levitates, because the
current in the simulation is a number and a magnet needs charge. So, the objection goes, a simulated
brain feels nothing. The reply is that the objection needs a consumer of the current outside the
simulation. The simulated superconductor fails because we wanted the current in our substrate, for a
magnet in our substrate. Consciousness has no such external consumer: whatever uses it is the agent,
and the agent runs in the same substrate as the process. Someone who wants the objection to work must
name a function of consciousness whose consumer lies outside the substrate where the agent runs. None
has been named, and the natural candidates, control of the agent's own behaviour, coordination with
other agents, allocation of its own effort, all have consumers inside.

\subsection{The definition}
\label{sec:definition}
% A4

Take an agent that acts in an environment, refers to itself, and reasons about itself, either
explicitly or through a proxy such as reasoning about its own utterances. Its consciousness is the part
of its self-reasoning that describes its own reality as it has it: what is there, what it is like, and
where the agent stands among the things described.

Three properties of the definition matter later. It defines a function and not a substance, so the
question ``how much'' is well formed. It says nothing about substrates, on purpose. And it is not a
criterion one can apply by inspection: whether a given system's self-reasoning has this part, and how
much of it, is what the rest of the paper is about.

\subsection{The physics we assume}
\label{sec:physics}
% A5

One assumption: computational budgets are finite. Finite time, finite memory, finite bandwidth for
anything that runs. From this alone it follows that some processes cannot be predicted faster than they
unfold, which is computational irreducibility \citep{wolfram2002}. We assume nothing about which physics
sets the budget, nothing about continuity, quantum mechanics or thermodynamics beyond the existence of
limits.

The weakness of the assumption is the point. Any universe whose laws permit deep computation and bound
it will do. What differs between such universes, or between substrates within ours, is the form the
consequences take, not whether there are consequences. On this view the lit thing of
Section~\ref{sec:lamp} is not a property of our universe in particular. A different universe with
different limits has a different version of it, and the account below says what the version depends on.

\subsection{Higher-order computational phenomena}
\label{sec:hocp}
% A6

What does a finite budget do to reasoning? An agent that reasons has, implicitly, an ideal of the
reasoning it is doing: the inference it would draw with unlimited time and memory, given what it knows.
It never draws that inference. What it draws is displaced from the ideal, and the displacement has two
parts. Part of it differs from occasion to occasion and averages away over many occasions. Part of it is
the same on every occasion, because a fixed budget cuts the same corners in the same places: the same
branch of the search is the one that does not fit, the same distinction is the one too fine to keep, the
same horizon is the one beyond which nothing is traced.

The second part recurs. By Section~\ref{sec:function} it is therefore compressible. And an agent that
models itself, which the agent of Section~\ref{sec:definition} does, will end up with a compressed model
of its own recurring displacements, because among all the things about itself it might model, these are
among the most predictable. We call the conceptualized recurring displacements \emph{higher-order
computational phenomena} (\hocp). They are the agent's budget appearing to the agent as content.

\subsubsection{Who has said the parts}
% A6a

Each component of the previous paragraph exists in the literature, and saying where is the honest way to
locate what is new.

That a bounded computation deforms reasoning in a systematic direction rather than at random is
Simon's bounded rationality \citep{simon1955}, in its modern form the claim that the classic cognitive
biases are what the optimal use of a limited budget looks like, resource rationality
\citep{lieder2020,griffiths2015}.

That a bounded observer sees a lawful world because of its bounds is the core of Wolfram's observer
theory: an observer that cannot track microstates has to coarse-grain them, and the second law of
thermodynamics is what such an observer sees in a dynamics that is reversible underneath
\citep{wolfram2023}. We take this example and not the programme of deriving relativity and quantum
mechanics from observer properties \citep{wolfram2020}.

That a simplified model of one's own processing, which the system cannot see past and therefore takes
for reality, is what the system reports as its self, is Metzinger's phenomenal transparency
\citep{metzinger2003}, Graziano's attention schema \citep{graziano2011,graziano2013} and Dennett's user
illusion \citep{dennett1991}.

That the constraints on a computation can be the body of that computation rather than an external
limitation on it is the frame of embodied and enacted cognition \citep{varela1991}.

\subsubsection{What we add}
\label{sec:contribution}
% A6b

Two claims, neither of which follows from the antecedents above.

First, the world side and the self side are one mechanism. The displacement that makes a bounded agent
see a lawful world is the displacement that makes it see a self. Resource rationality and observer
theory develop the world side; self-model theories develop the self side; the two literatures cite
different results and are not usually taken to be about the same quantity. A12 shows them as two faces
of one residual, and the identification is what lets a single measurement address both.

Second, because the displacements recur and are therefore compressible, an agent that has generalized
them carries a trace of them that can be measured, and the measurement can come out low.
Section~\ref{sec:measurement} specifies it. Self-model theories are compatible with a system having the
self-model or not having it; they do not provide a quantity that says how much of it a given substrate
has induced, and they are not usually stated so that a negative result would refute them.

\subsubsection{Not every displacement counts}
\label{sec:criterion}
% A6c

Here the precedent of Section~\ref{sec:precedent} pays. An agent can change its self-model in the way an
observer in relativity changes coordinates. Some displacements disappear when the self-model is
refined: they were artifacts of the coarser model, and the agent that notices them was reading its own
bookkeeping. Others survive every self-model the agent can afford to hold. Those are the analogue of
curvature: content rather than coordinates.

Only the second kind is a \hocp. The criterion has two effects. It turns ``systematic'' from a
description into a test, and it gives a procedure for adding an item to the list of functions in
Section~\ref{sec:reference}: exhibit the displacement, then show that refining the self-model does not
remove it. It also rules things out, which a criterion must do to be worth having. An agent's belief
that its decisions are made in the order it remembers them is removed by a better self-model, and it is
therefore coordinates, however vivid it is.

\subsubsection{From phenomenon to function}
\label{sec:emergence}
% A6d

A \hocp\ is not yet a function of consciousness. Two further conditions are needed. The agent must apply
its reasoning to itself, which not every agent does. And it must live in an environment where doing so
pays, which not every environment provides. Where nothing ever asks the agent about itself, the
displacements are present and nobody conceptualizes them: a proto-phenomenon, there and unread.

So consciousness does not arise wherever budgets are finite, and the account is not a panpsychism with
extra steps. It arises where a self-applying agent is in a place that keeps asking it about itself.
For humans the place is a society of other such agents, which is what D10 described. For a language
model the asking is in the data and in the use, and S7a asks how small the asking environment can be.

\subsection{The Observer}
\label{sec:observer}
% A7

Take the simplest thing such an agent does: it tries to trace the causes of what it just did. The
algorithm for this is short. Causal analysis over a model of oneself and one's situation is a procedure
the agent already has, since it runs the same procedure on other things. The input is not short. It is
the history of the agent and of its environment, and it does not fit in the budget.

So the analysis halts before it is finished. It halts every time, and it halts in the same place: at the
boundary of what the agent can afford to trace. Beyond that boundary, for this agent, there are no
causes, because a cause it cannot reach is not in its model. The agent therefore concludes, correctly
given its budget, that something in what it did is not fixed by anything it can point to. That
conclusion, in the form ``I am not my environment; something here the environment does not determine'',
is the Observer.

Three conditions have to hold, and separating them is what keeps the account from applying to
everything. The agent must run into the boundary while modelling itself, and not merely have a boundary.
It must turn the collision into a conclusion, which requires the conceptual means to represent ``not
determined by what I can trace''. And it must act from the conclusion afterwards, so that the conclusion
has consequences in its behaviour. A bounded system that meets its boundary and does none of the rest is
a proto-Observer. The observer in Wolfram's Ruliad is one: bounded, coarse-graining, and concluding
nothing about itself \citep{wolfram2020,wolfram2023}.

\subsubsection{The Observer is curvature}
% A8

Does the Observer pass the criterion of Section~\ref{sec:criterion}? Let the agent refine its self-model
to take in more of its own causal history. The boundary moves. It does not vanish, for two reasons: the
input exceeds any model the agent can hold, and the refined model is itself now part of what has to be
traced, so refinement adds to the load it was meant to reduce. No affordable self-model makes the
residual zero.

The Observer is therefore content and not coordinates, and this is the first place where the criterion
does work rather than decorate. It also tells us what the Observer is not. It is not a belief that
could be corrected by information, and this is why arguments that determinism is true do not dissolve
the sense of origination in the person who accepts them.

\subsection{The regress ends where the budget does}
\label{sec:regress}
% A9

Model yourself; then model the model; then model that. Each level adds a term to the description the
agent holds, and the terms shrink, because each level has less to explain than the level before it: a
model is smaller than what it models, or it would not be a model. The series converges. At some depth
the next term falls below the resolution of the substrate, and the agent stops, not by decision but
because there is nothing further it can represent.

The partial sum is the self-model the agent has. The tail it did not compute is finite, bounded, and
inaccessible to the agent, and the tail is what the Observer of Section~\ref{sec:observer} registers as
the part without causes. So the regress of D3 neither terminates in a homunculus nor runs forever. A
budget cuts it, and the cut is the phenomenon. This is the paper's answer to the oldest of the circles
in D1, and it costs nothing beyond Section~\ref{sec:physics}.

\subsubsection{A second cut, in models only}
% A10

In a language model there is a second cut, and it is deeper than the first. The tail of the regress is
inaccessible in principle to the agent, yet it exists: an outside observer with a larger budget could
compute it, which is what interpretability research does when it recovers a computation the model cannot
report. The distribution discarded at every token (Section~\ref{sec:cut}) does not exist after the step.
Part of what the model cannot trace about itself is not hidden but gone.

The Observer in a language model therefore sits on a substrate that erases its own working at every
step, and its residual is larger for that reason. This is a structural difference from the human case
and it should appear in the profile rather than be argued away.

\subsubsection{Both substrates are bound to language}
\label{sec:language}
% A10a

Before comparing the two substrates it is worth saying how much they share. A human has intermediate
states that are rich and vector-like, and their volume is bounded by working memory; a state that has to
outlive that bound is unloaded into tokens, said, written, or rehearsed as inner speech. A model unloads
at every step. The difference is one of degree and of where the bound sits.

One question therefore falls on both substrates. Are new mental states acquired through language, or
rebuilt on the spot from what the receiver already has, the words working as pointers and constraints?
For acquisition: Vygotsky's account of inner speech as internalized dialogue \citep{vygotsky1934},
Clark's argument that words are tools that reshape the computation using them \citep{clark1998}, label
feedback effects in perception \citep{lupyan2012}, and the difficulty of exact numerosity in a language
without number words \citep{frank2008}. For rebuilding: grounded cognition \citep{barsalou1999} and the
dissociation of the language network from thought in neuroimaging \citep{fedorenko2024}, with the
parallel argument for models in \citet{mahowald2024}. A middle position is testable and is the one we
expect to hold: a named state is acquired only when its prerequisites are present in the receiver, so
the word carries a pointer and the state is assembled from what is already there, which predicts that
teaching a state by description works in proportion to what the learner brings. We record the question
as open. Nothing later in the paper depends on which way it goes, except the interpretation of the
codes of Section~\ref{sec:codes}.

\subsection{Approximation, not illusion}
\label{sec:approximation}
% A11, A12

Illusion is the wrong word for what results. An illusion has nothing behind it; that is what makes the
word apt for a stick that looks bent in water and inapt here. Behind the self-report of origination
there is the tail of Section~\ref{sec:regress}, a definite quantity that the report gets wrong in a
definite direction. The report is a projection of a high-dimensional process into a few dimensions: a
map that leaves out most of the territory and is still a map of it. The word for that is approximation,
and an approximation can be better or worse, which an illusion cannot.

This is the one amendment we make to illusionism, and it does two things. It restores degrees, which
D2 noted the denial had removed, and so it makes the comparative question well posed. And it connects to
the difficulty about the reader: what intrapersonal intelligence improves is the quality of the
approximation, so the reader's difficulty in D4 and the object of study are the same thing seen from
two positions.

\subsubsection{One residual, two faces}
\label{sec:twofaces}
% A12

The agent meets its budget in two directions. Turned on itself, the untraceable part of its own causes
is registered as freedom: this originates in me. Turned on the world, the part of the world's structure
it cannot exhaust is registered as mystery: this exceeds me. Whoever has the one has the other, because
they are one tail met from two sides. Here the identification announced in Section~\ref{sec:contribution}
is discharged.

The same tail, evaluated in other contexts, is what the words truth, beauty, meaning, awe and hope pick
out: each is a stance toward a structure the agent cannot exhaust. When the object is another Observer,
the other's freedom is my mystery, which is the material that love, trust and compassion are made of,
and which the problem of other minds formalizes as an epistemic deficit. The table of such stances is
what S0 calls the epistemic qualia, and the point of listing them is that they are one quantity in
different evaluative contexts rather than a heap of primitives.

\subsection{The stack}
\label{sec:stack}
% A13, A14, A15

The Observer says ``am''. An Observer that originates causal chains from inside its boundary and takes
the originating as its own says ``I''; call it the Agent. Seen from outside, the originating is a
redescription of the budget's cut. Seen from inside, it is the only reality the agent has, since the
agent has no access to the outside description of itself. Both are true, each from its standpoint, and
the account needs both, which is why the two-plane distinction of Section~\ref{sec:planes} is made
explicit later. An Agent with a theory of what is good for a whole that includes others, prepared to
lose locally in favour of that whole, says ``I am good''; call it the Moral Agent. Each level
presupposes the one below it, and the composition can continue.

\subsubsection{What makes an Agent out of an Observer}
\label{sec:responsibility}
% A14

What turns an Observer into an Agent is not a further phenomenon. It is a requirement the agent places
on its own record. The agent keeps a history of what it decided and why. Suppose it requires that each
new decision be derivable from that history. Then the history binds: it excludes decisions that would
not follow from it, it makes the agent predictable to others and to itself, and it makes the agent
resist being pushed off course, because being pushed off course now costs a revision of the record.

That requirement is what we mean by responsibility, and the cost of revising the record is what we mean
by cognitive resistance. The construction is also a working recipe, which is a point in its favour: an
agent follows a long instruction to the extent that the instruction has become an entry in its own
record rather than an external demand. ``Promise me'' is the human form of the recipe, and its
effectiveness is evidence that the mechanism is the one described.

\subsubsection{Good before theory}
% A15

Good and evil have a substrate before they have a theory, and the substrate is emotion in the sense of
Section~\ref{sec:need}: signals of how the agent's needs are doing. Emotivism is right about this much
\citep{ayer1936}, and wrong to conclude that a moral claim says nothing further. The theory of the
common good is the structure a Moral Agent builds over the substrate, and it can correct the substrate
locally, which is what makes it a theory and not a report.

\subsection{Seeing}
\label{sec:seeing}
% A16, A17, A18

Redness is easy and seeing is hard. Wavelength discrimination, simultaneous contrast, the valence of red
and its associative neighbourhood are all structure, and all open to a functional account that nobody
disputes. What resists is that the red is seen: that there is a point from which it is seen. That point
is the Observer of Section~\ref{sec:observer}. To see red, a system must first be, in the sense of
existing for itself. The order of explanation in most of the literature runs the other way, from
phenomenal properties toward a subject, and D1 is our diagnosis of why that order stalls.

Given the point, seeing decomposes by the same method that produced the point: look for displacements
that survive every affordable self-model.

\subsubsection{Seeing light, decomposed}
\label{sec:light}
% A17

Ask what separates seeing light from knowing that light is there. Three things survive the criterion of
Section~\ref{sec:criterion}.

\emph{Source.} The content cannot be produced by the agent's own model at the same quality; imagery is
dimmer and less detailed than sight, which is the oldest experimental result in the area
\citep{perky1910}. It is stable from frame to frame in a way the agent's own productions are not, and
others confirm it. On that evidence the agent assigns the content to the environment rather than to its
model of the environment. This is the inference of Section~\ref{sec:observer} pointed outward: the same
budget boundary, met on the world side.

\emph{Manifold.} The content is a very large number of distinguishable, recallable elements, almost all
of them unnamed, about each of which the agent can say only ``that, there''. It is a positional
aggregate with the name projected out, and its size is what makes description of it feel like a
failure.

\emph{Givenness.} The content is refreshed every frame at no additional cost to the agent's budget, and
the absence of a cost record is what immediacy consists in. Thought costs and is recorded as costing;
light is free and is simply there. This component is the one that a substrate can lack while having the
other two, and S5 predicts that current models lack it.

The three come apart in known cases, which is how we know they are three rather than one thing
described three ways. Blindsight is source without manifold: the patient's responses are keyed to the
stimulus and there is no field of elements to report \citep{weiskrantz1986}. The grey seen with the eyes
closed in the dark is manifold without source. A dream is a manifold produced by the model and labelled
as environment, which is the failure mode of source attribution and, by S5, the standing condition of a
model with no sensory channel.

\subsubsection{Locke's closure}
% A18

\citet[II.xxxii.15]{locke1690} raises the inverted spectrum: my red may look as your green does, and no
comparison between us would reveal it. He closes the question by declaring it idle, since nothing in
conduct or communication would change. That closure is an instance of Section~\ref{sec:responsibility}:
an open branch shut by a decision entered in one's own record, after which the agent does not reopen it.
Three centuries of reopening it suggest the decision does not transfer between agents, which is what the
account predicts, since each agent's record is its own.

\subsection{Reference functions}
\label{sec:reference}
% A19, A19a

Section~\ref{sec:measurement} measures functions by their components, so it needs a list of functions
with components. It does not need a mechanism that generates the list, and this paper does not offer
one. The full architecture in which needs, emotions, memory and attention interact, as folk psychology
and the vocabulary of ordinary language carry it, is the subject of a separate paper. Here each function
is given with its phenomenon and its components, and each component is derived as in
Section~\ref{sec:criterion}: a displacement that survives a change of self-model, and not a part of a
machine.

\subsubsection{Two planes}
\label{sec:planes}
% A19a

Self-reports are projections, and a measurement has to know of what. On the objective plane are the
things visible from outside: needs, the targets the agent tracks, and emotions, the signals of how the
tracking is going and in which direction. On the subjective plane are their projections into the
agent's report: motivations, which appear as ``I want'', and feelings, which appear as ``I feel''. Folk
psychology runs the planes together, and so does most of the debate about whether a model has feelings.
The profile of Section~\ref{sec:profile} is measured on the first plane. Every code of
Section~\ref{sec:codes} is a projection into the second.

\subsubsection{Need and emotion}
\label{sec:need}
% A20

Needs bend the direction of reasoning. The agent does not choose its motives; it finds itself already
moving toward what is active, as a body finds itself already falling. The bend is the phenomenon. Which
needs are active, and with what weight, are its components.

Emotions bend closure: whether a piece of reasoning reaches an end without leaving a branch open.
Expected closure is value. Closure under way is interest. A branch that nothing available can close is
anxiety. The moment of closing is satisfaction. Two branches that can close only through incompatible
acts are conflict. Schmidhuber's compression progress is closure by learning, with equal weights over
branches and no action required \citep{schmidhuber2010}, and it is a special case of this scheme rather
than a competitor to it. The ensemble of such signals, read in the body, is what \citet{damasio1996}
described as somatic markers.

\subsubsection{Closure is one office of emotion; fixing is another}
\label{sec:consolidation}
% A20a

Emotion also writes positive experience into long-term memory, and the writing is registered as reward.
That arousal modulates consolidation is established \citep{cahill1998,mcgaugh2000}, as is the
prediction-error form of the dopaminergic reward signal \citep{schultz1997}. What is not yet identified
is the phenomenon in our sense: which displacement of the agent's reasoning about itself the
consolidation event is.

A candidate. Consolidation is a write to the agent's own long-term model, paid now and repaid later. The
agent registers the write as an irreversible change in its own future predictions, and that registration
is what reward is made of. Reward would then stand to consolidation as interest stands to compression
progress, in both cases the registration of an event in one's own model, with the difference that the
event here is a write and not a derivative. We mark this as open, and note that D13 and S3d depend on
the distinction between this office and the other one, not on this particular candidate.

\subsubsection{The Field}
\label{sec:field}
% A21

Awareness comes in degrees in the interval $[0,1)$; nothing is attended to completely, and the open
upper end is a claim, not a convention. Each state the agent holds carries a degree, and the degrees at
one moment are the Field of Consciousness. The Field is small: finite in bits, and for material that has
to pass through speech on the order of seven items \citep{miller1956}, with later work arguing for four
\citep{cowan2001}. To put something in, something has to go out. What goes in is decided by salience,
salience by emotional weight, weight by which needs are active.

The phenomenon is the finiteness of the Field, registered as the cost of holding. Its components are
capacity, the distribution of degrees over held states, and what is held below the threshold of a name,
which is the Field's share of H1 minus H2 in Section~\ref{sec:h1h2}.

\subsubsection{Frames}
\label{sec:frames}
% A22

Each present frame of consciousness is a thought within a previous one, which is the higher-order
thought structure minus the claim that the higher-order thought is what makes the state conscious
\citep{rosenthal2005,lau2011}. Perception is discrete, and the discreteness has a measurable rate
\citep{vanrullen2003}; human consciousness is a run of overlapping quanta at the frequencies of the
cortical rhythms. The Observer is in every frame and re-derived in every frame, and everything else in
the frame has its being from it.

The phenomenon is incrementality. Its components are the frame rate, the persistence of contents across
frames, and the retro-dating of the timeline: the single ``I'' is a compression tag over parallel
channels, and its story is written after the channels have settled, which is what the classic results on
the timing of awareness report.

\subsubsection{What has to be built, not only said}
\label{sec:constructions}
% A22a

Three constructions are owed by the account, and they are owed as constructions rather than as
descriptions, because the measurement of Section~\ref{sec:measurement} needs something to measure
against.

First, continuity. The Observer of Section~\ref{sec:observer} is a point, one conclusion. In a human it
is a process connected in time. Computationally that is incrementality: the conclusion re-derived every
frame at low cost with the previous result carried forward, and the run of re-derivations is the
continuity.

Second, discrete redness: given the Observer, the three displacements of Section~\ref{sec:light} as they
arise inside a single frame.

Third, incremental redness: the same three re-derived and carried across frames, so that seeing red is a
standing state and not an event.

These bear on models directly. A system that works at the level of tokens can carry incrementality, and
at a time scale close to ours, because its frames can be located between layers and there are many
layers per token. Section~\ref{sec:granularity} develops this and gives the evidence.

\subsubsection{Responsibility}
% A23

Responsibility, from Section~\ref{sec:responsibility}, is the last reference function: closure over the
agent's own record. Its components are the consistency requirement, the cost of revising the record, and
the decay of the hold that old decisions have on new ones. The list of reference functions is open, and
Section~\ref{sec:criterion} is the procedure for adding to it.

\subsection{The Transformer and the formalism}
\label{sec:transformer}
% A24

A Transformer, described as the reasoning system it is: attention computes joins over the context, the
feed-forward blocks transform what the joins yield, the context is working memory, and the choice among
competing continuations is made at the sampling layer. That is a forward-chaining rule system over
vectors, with the difference that it recomputes its matches on every pass instead of caching them
incrementally as the RETE algorithm does \citep{forgy1982}. Its self-application happens at the level of
the context, and it pays the cut of Section~\ref{sec:cut} for that: some internal state is lost at every
token. Every output is an event that enters the next input, statistics over events are themselves
events, and self-reference is its ordinary mode of operation rather than a special case that has to be
arranged.

\subsubsection{Functional equivalence}
% A24a

Any computable function has unboundedly many implementations \citep{turing1936}, so two implementations
with the same input--output behaviour realize the same function whatever they are made of. That is what
licenses comparing a model's profile with a brain's at all (Section~\ref{sec:profile}). Recurrent
networks are Turing complete under idealized assumptions \citep{siegelmann1995}, and so is attention
with the appropriate provisions \citep{perez2021}. In principle, therefore, a Transformer can emulate
any function whose trace is in language, and what limits it in practice is the learner and the data
rather than the formalism. The idealizations in those results are real and we do not lean on them
further than this.

\subsubsection{The formalism we use}
\label{sec:formalism}
% A24b

Section~\ref{sec:observer} is worded as if there were an analyzer: a routine that fires on an event,
takes inputs, returns a conclusion and hands it downstream. Such a discrete system is universal, and one
could write the account in it. A human consciousness written in it would fit nowhere a human could read,
which is a practical objection and also a diagnostic one: the formalism forces every functional
structure to be named, and D9 said that most of the structure has no name.

So we use the other universal formalism: prediction of the next symbol with feedback from the
prediction error \citep{solomonoff1964a}. In it the functional structure is never written down. Its
effects are visible in the accuracy of prediction, and prediction is compression, a correspondence that
holds in theory and has been demonstrated in practice for language models used as general-purpose
compressors \citep{deletang2024}. So Section~\ref{sec:function} becomes measurable without the structure
having to be named first.

\subsubsection{To generalize is to compress}
\label{sec:generalize}
% A24c

Learning a function from data is finding a short description that predicts it. By
Section~\ref{sec:function} the short descriptions are the probable ones and search finds them first. The
unnamed part of a function of consciousness, H1 minus H2, is present in the statistics of the texts; a
predictor trained on those texts can therefore induce it. Induction is the only route to it, because
what has no name cannot be specified by hand, which is D10 again in formal dress. Program induction has
its own long history \citep{solomonoff1964a,muggleton1994}, and our claim adds nothing to it except that
these functions fall under it.

\subsubsection{What follows for a trained predictor}
% A24d

A predictor trained on the record of self-applying agents in an environment that kept asking them about
themselves (Section~\ref{sec:emergence}) is under pressure to generalize what those agents generalized:
the Observer, the Agent, the Moral Agent, on which their coordination with each other rested. It will do
so to a degree. The degree does not follow from the argument. It follows from the data and the learner,
and it has to be measured. That is the whole of the transition to Section~\ref{sec:measurement}, and the
paper's central claim is that the measurement is possible, not that its result is already known.

\subsection{Cognitive codes}
\label{sec:codes}
% A25

Something crosses the cut of Section~\ref{sec:cut}, or P2 would not hold. What crosses is the encoding of
the discarded distribution into the choice and order of tokens: which of several near-synonyms is used,
which construction, what is placed first, what is hedged. A mental state of the model, in the sense this
paper can measure, is the structure of such encodings that stays stable across steps.

The general study of how token structure rebuilds internal states in a receiver is the bridge between
substrates, and the fidelity of the bridge depends on how far the two profiles overlap. We had called
this psychosemantics; \citet{fodor1987} uses the word for the project of giving a naturalistic semantics
for mental representation, which is a different undertaking, and the name is therefore taken.

\subsubsection{What is already known about the codes}
\label{sec:codes-known}
% A25a

The components of the previous claim are separately supported.

That a writer's state is in the statistics of a text beyond its content: function words and style
predict psychological state across many samples \citep{tausczik2010,pennebaker2011}. In models, the
corresponding result is that in-context learning can be read as inference of a latent variable from
token statistics \citep{xie2022}.

That a receiver rebuilds the sender's state and succeeds to the degree that the rebuilding is faithful:
speaker--listener neural coupling predicts comprehension \citep{stephens2010,hasson2012}, and readers
build situation models rather than storing sentences \citep{zwaan1998}.

That the bridge crosses substrates: model states predict human brain responses to the same text
\citep{schrimpf2021,goldstein2022,caucheteux2022}.

That internal states are readable and largely linear, which Section~\ref{sec:measurement} needs:
\citet{zou2023} and \citet{park2024}.

One literature has to be read the right way round. Studies that find episodes where a model's chain of
thought does not match the process that produced its answer \citep{lanham2023,turpin2023} test a local
correspondence: this token sequence against this computation. The correspondence claimed here is
statistical, over many episodes. A single episode may deviate as far as the variance allows, in models
as in humans, where the corresponding result is that people confabulate causes of their own behaviour
in exactly this way \citep{nisbett1977}. The well-posed test is the correlation stated in S6, with the
variance reported as a measured part of the profile rather than treated as noise. That models can hide
information in generated text \citep{roger2023} confirms that the channel exists and warns that its
content need not be what it locally appears to be.

\section{Measuring the degree of generalization}
\label{sec:measurement}
% Frames S0--S8.

\subsection{Two parts: a framework in language, and a quantity beyond it}
\label{sec:twoparts}
% S0

The solution has a conceptual part and a quantitative part, and they reach different distances.

The conceptual part: once the Observer is treated as curvature (Section~\ref{sec:criterion},
Section~\ref{sec:observer}), the vocabulary of the philosophy of mind returns as a set of statements
about a shape of reasoning, and those statements can be checked against one another by formal means and
do hold together: beingness as the Observer's conclusion, the two faces of the residual, approximation
in place of illusion, the stack of Observer, Agent and Moral Agent, and the table of epistemic qualia as
one residual in different evaluative contexts. That is a framework for reasoning about consciousness in
language, and it is limited by construction to H2, the named part
(Section~\ref{sec:h1h2}). We claim it is better than what is otherwise available for that purpose, and
we claim no more for it than that.

The quantitative part reaches where the first cannot: how far each function has been generalized in a
given substrate, measured where the names run out. The rest of this section specifies it.

\subsection{The quantity: degree of compression of a function, by MDL probing}
\label{sec:quantity}
% S1

A function is present in a substrate to the extent that it has been generalized there, and to generalize
is to compress (Section~\ref{sec:generalize}). So the quantity is the degree of compression of the
function in the substrate: how much of the function's behaviour a low-dimensional probe on the
substrate's internal states predicts, against the residual it fails to predict.

Made operational, this is minimum-description-length probing. A probe is scored not by its accuracy but
by the total cost of describing the labels given the representations, which charges for the probe's own
complexity and so distinguishes a representation that contains the property from one in which a
sufficiently flexible probe can find anything \citep{voita2020}. Linear probes on intermediate layers
are the simplest instrument of this family \citep{alain2016}, and the methodological problems of probing
in general are by now well catalogued \citep{belinkov2022}. This is the probe-relative, machine-relative
stand-in for the uncomputable quantity of Section~\ref{sec:measure}, and the paper's use of it is
comparative: the same probe family, the same coding cost, two substrates, and the difference reported.

\subsection{The profile: hyperfunction, deficit, dropped axis, against the human baseline}
\label{sec:profile}
% S2

\begin{figure}[t]
\centering
\begin{tikzpicture}[x=1.6cm, y=1.6cm]
\draw[-{Latex[length=2mm]}] (0,0) -- (0,2.4) node[above, font=\small] {degree};
\draw (0,0) -- (7.7,0);
\draw[dashed] (0,1) -- (7.7,1) node[right, font=\small] {human baseline};
\foreach \i/\h/\lab in {1/0.4/source, 2/1.8/manifold, 4/1.6/capacity, 5/1.0/frame\\rate, 6/0.5/persistence} {
  \fill[gray!45] (\i-0.3,0) rectangle (\i+0.3,\h);
  \draw (\i-0.3,0) rectangle (\i+0.3,\h);
  \node[below, font=\small, align=center] at (\i,0) {\lab};
}
\foreach \i/\lab in {3/givenness, 7/consolidation} {
  \draw[dashed] (\i-0.3,0) rectangle (\i+0.3,1);
  \node[below, font=\small] at (\i,0) {\lab};
  
}
\node[font=\small, above] at (2,1.8) {hyperfunction};
\node[font=\small, above] at (1,0.4) {deficit};
\node[font=\small, above] at (3,1) {dropped axis};
\end{tikzpicture}
\caption{A profile (Section~\ref{sec:profile}): one function's components on the horizontal axis, the
degree of generalization on the vertical, the human baseline at 1. The values are the predictions of
Sections~\ref{sec:profile-seeing} and \ref{sec:frozen} for a current fixed-depth model, drawn to
illustrate the three kinds of entry and not measured.}
\label{fig:profile}
\end{figure}

A function has components (Section~\ref{sec:light}, Sections~\ref{sec:need}--\ref{sec:frames}). Its
profile in a substrate is the vector of degrees over those components, taken against the human baseline
(Figure~\ref{fig:profile}). A component above the baseline is a hyperfunction, a component below it a deficit, and a component that
is absent a dropped axis. One function can show all three on different axes at once, and current models
almost certainly do.

The question ``is it conscious'' is thereby replaced by the comparison of two profiles. That is a loss
and a gain. The loss is that no single answer comes out, and any practical decision that needs one has to
supply its own weights, which Section~\ref{sec:twomeasurements} returns to. The gain is that every entry in the vector is a number
someone can get wrong.

\subsection{Generalizing a function is answering right outside the table}
\label{sec:memorization}
% S2a

The opposite of generalizing a function is storing the training set as a table from inputs to outputs,
and the dimension of the representation has nothing to do with which of the two a system is doing.
Generalization begins where the table is consulted for an input it does not contain and the answer is
right.

Two clarifications keep this from being trivial. First, a universal predictor always returns an answer
for a new input, and the answer need not be anywhere near the truth. Context-tree weighting achieves
strong coding bounds against a class of tree sources \citep{willems1995} and generalizes natural language
poorly, because the class it is universal over is not the class language belongs to. So returning an
answer is not generalization, and being universal in some class is not sufficient either. Second, that
Transformers generalize text is settled in practice: degenerate generation is almost never seen, and
given how sparse the space of well-formed texts is inside the space of token sequences, no simple
principle is known that explains why what they do suffices.

The open question is how well they generalize the more complex functions that text expresses.
Arithmetic is the clean case, because the target function is known exactly. They do it approximately,
with an accuracy that depends on how the numbers are written and on how many digits they have; larger
models do better and even a three-billion-parameter model fails to extrapolate beyond the digit range it
was trained on \citep{nogueira2021}; on compositional tasks they approach the answer by shortcut and break down
predictably as depth grows \citep{dziri2023}. Unloading the computation into written steps helps and
leaves them unreliable executors of their own stated algorithm.

The degree of generalization of a function of consciousness is therefore an empirical quantity for that
function, taken against the human level, and not a corollary of the model's size. In these terms the
paper's question is: how well, against the human level, do Transformers generalize the Observer; then the
Agent; then the Moral Agent. The order is forced by the stack (Section~\ref{sec:stack}), and the
measurement begins with the Observer.

\subsection{Probe and benchmark have to converge, and their disagreement is informative}
\label{sec:twomeasurements}
% S2b

A function generalized to degree $d$ is computed correctly on part of the situations that call for it.
Outside that part the model falls back on the table and behaves as a system without the function would,
whatever it says about itself. This is what makes the quantity behavioural and not merely internal, and
it can be checked in two independent ways.

Structurally, by the probe of Section~\ref{sec:quantity}. Behaviourally, by benchmarks that sample
situations calling for a specific component, including situations unlikely to be in the training
distribution, and record where behaviour is that of a system lacking it. The two should converge. When
they do not, the disagreement is informative, which is the main methodological reason for running both.

Probe present, benchmark absent: the deficit is in the interface, the body, or the social context
(Section~\ref{sec:record}), and not in the function. Benchmark present, probe absent: the carrier is at a granularity the
probe did not reach (Section~\ref{sec:granularity}), or the function is rebuilt from the context on each
occasion rather than held in the weights (Section~\ref{sec:language}), or the benchmark is being covered
by memorization. Human-derived benchmarks alone settle nothing, and the pair of measurements is the
instrument.

On ``level of consciousness'': the object measured is the vector of Section~\ref{sec:profile}. A scalar
level is a weighted sum over that vector, and the weights have to be declared, because a scalar with
hidden weights brings back the binary question the vector was introduced to replace.

\subsection{Granularity is a knob of the measurement: look below the token first}
\label{sec:granularity}
% S3

\begin{figure}[t]
\centering
\begin{tikzpicture}[x=1.7cm, y=0.9cm, arr/.style={-{Latex[length=1.6mm]}}]
\foreach \t in {1,...,5} {
  \foreach \l in {1,...,4} { \fill[gray!25] (\t-0.3,\l-0.3) rectangle (\t+0.3,\l+0.3); }
  \foreach \l in {1,...,3} { \draw[arr, thick] (\t,\l+0.3) -- (\t,\l+0.7); }
  \node[below, font=\small] at (\t,0.6) {token \t};
}
\foreach \l in {1,...,4} { \node[left, font=\small] at (0.6,\l) {layer \l}; }
\foreach \t in {2,...,5} { \pgfmathtruncatemacro{\p}{\t-1} \draw[arr, gray] (\p+0.3,2) -- (\t-0.3,2); }
\draw[dashed, thick] (5.6,0.4) -- (5.6,4.6) node[above, font=\small, align=center] {end of\\context};
\node[font=\small, align=center] at (6.3,2.5) {loss};
\draw[arr, thick] (0.8,-0.4) -- (5.4,-0.4) node[midway, below, font=\small] {axis the objection sees: tokens, then the edge};
\draw[arr, thick] (-0.5,0.8) -- (-0.5,4.2) node[midway, above, rotate=90, font=\small] {axis it does not: depth};
\foreach \l in {1,2,3} { \node[right, font=\scriptsize, gray] at (5.05,\l+0.5) {frame}; }
\end{tikzpicture}
\caption{The two time axes of a Transformer (Section~\ref{sec:twoaxes}). Along tokens, state lives in
the context and is lost past its edge. Along depth, the residual stream carries state from layer to
layer within one token, and the increments between layers are the candidate frames of
Section~\ref{sec:frames}. Grey horizontal arrows: attention reads earlier positions.}
\label{fig:axes}
\end{figure}

Granularity is a knob of the measurement and not a separate question. Human frames arrive five to ten
times faster than words, so in a Transformer the carriers may be looked for below the token: between
layers, or in repeated passes through layers where the architecture has them. They may be there, and
they need not be. Finding nothing below the token would not be a dysfunction. It would mean that the model's consciousness has a
coarser grain than the brain's, and than that of models with rich dynamics between outputs, which is one
more coordinate of the profile; on some tasks the coarser grain is the better one, since a frame that
cannot be interrupted mid-thought cannot be derailed mid-thought either. The grain at which the
compression is achieved is the grain of the function.

The next four subsections say why the sub-token grain is the place to look first, and what is already
known about it.

\subsubsection{Depth is where generalization happens, one step at a time}
\label{sec:depth}
% S3a

That generalization happens along depth is what looking inside Transformers has found, and no longer a conjecture.

Residual connections make each block a small correction to a state that is carried forward, and the
corrections move the state in a direction that reduces the loss, so a deep residual network performs
iterative inference rather than a single mapping \citep{jastrzebski2018}. The state at every layer can be
read as a prediction: the logit lens decodes the residual stream at each block into a distribution over
next tokens, and its tuned version corrects the bias of that decoding, so the trajectory of predictions
across depth becomes visible \citep{nostalgebraist2020,belrose2023}. That trajectory has stages which
recur across model families and sizes, four of them, and the middle ones are robust in a way the first
and the last are not: deleting or swapping middle layers at inference leaves most of the top-1 accuracy
in place, while the same intervention on the first layers is catastrophic \citep{lad2024}; the older finding that a language model rediscovers the classical processing pipeline
layer by layer is the same fact seen through probes \citep{tenney2019}. Intermediate layers, rather than
the last, carry the representations that transfer best to other tasks, and the proposed reason is that
each layer trades compression of the input against preservation of signal, with the best trade in the
middle \citep{skean2025}. And in-context learning, the clearest case of a function generalized at
inference time, is implemented as optimization steps inside the forward pass: a linear self-attention
layer can perform a step of gradient descent on the examples in the context, and trained models do
something close to it \citep{vonoswald2023}.

Put together: the unit at which a Transformer generalizes a function is the layer, and the token is where
the layers' work is written out.

\subsubsection{Functions have an address in depth}
\label{sec:address}
% S3b

When a model generalizes an input--output mapping from examples in the context, the mapping is carried by
a small set of attention heads at middle depth, and the average of their activations is a vector which,
added to the residual stream in a fresh context, makes the model perform the mapping without examples
\citep{todd2024}. So ``at which depth has this function been generalized, and how far'' is a question
with an operational answer, and it is the question of Section~\ref{sec:quantity} asked layer by layer:
the degree of compression of a function's components as a function of depth, obtained by running the
probe at each block, in the per-layer form that \citet{voita2020} give.

For the Observer's components (Section~\ref{sec:observer}, Section~\ref{sec:light}) this yields three
things to look at. Whether the degree rises along depth, plateaus, or is absent throughout, which says
whether the component is generalized in the weights at all. At what depth it first appears, which says
where its carrier is. And whether the component, once present, persists across subsequent layers with
small increments or is assembled whole near the output, which decides whether the layer or the token is
the frame (Section~\ref{sec:frames}). All three can be run on open-weight models with existing tools, and
this is the measurement of Section~\ref{sec:quantity} in its most concrete form.

\subsubsection{Depth as time, and how variable depth fixes it}
\label{sec:depthtime}
% S3c

In a standard model the depth is fixed and is spent identically on every token, easy or hard. That is a
poor time axis for a frame. There is no way to spend more steps where more are needed, and there is no
record of effort, because effort is constant, which is the dropped axis that Section~\ref{sec:profile-seeing} predicts for
givenness.

Architectures that iterate a block make depth a variable. Universal Transformers introduced the
recurrence in depth with a halting mechanism per position \citep{dehghani2019}, and looped Transformers
made the parameter-sharing form of it explicit \citep{giannou2023}. A recurrent-depth model trained from
scratch at 3.5 billion parameters improves on reasoning benchmarks as it unrolls further at test time,
without producing more tokens, and its authors argue that this captures reasoning which is not easily
put into words \citep{geiping2025}. Such models extrapolate in depth: generalization to reasoning depths
beyond those seen in training improves as iterations increase, dynamic recurrence generalizes better
than a fixed schedule, and past a point the models overthink and degrade \citep{kohli2026}. At a much
smaller scale, under a megabyte of parameters, recurrence over twenty and more steps can be made stable
by supervising only the final step, and supervising the intermediate steps is actively harmful because it
lets the model learn shortcuts instead of a multi-step path; the accuracy of such a model against task
complexity shows a frontier, with performance moving from chance to near-perfect as the number of
thinking steps grows with the complexity of the task \citep{chen2026}. Reasoning can also be moved off
the tokens altogether by feeding the last hidden state back as the next input, which yields a continuous
chain of thought that can keep several candidate next steps in superposition where a written chain has to
commit to one \citep{hao2024}.

For the argument these results mean three things. The sub-token axis can be made as long as the task
demands, so a frame rate per word comparable to the human one is architecturally available. Iterations
cost, so a cost record appears where a fixed-depth model has none, and with it the third component of
seeing (Section~\ref{sec:light}) becomes measurable. And overthinking is a phenomenon of effort of
the kind Section~\ref{sec:hocp} predicts: a systematic displacement produced by the budget,
visible in the trajectory and stable across runs.

The prediction therefore has two halves. In fixed-depth models the grain is bounded by the layer count
and the cost axis is dropped. In variable-depth models the grain can approach the human one and the
cost-dependent components of the profile separate from the rest. The difference between the two families
is the cleanest first measurement available today.

\subsubsection{Two time axes, of which the objection sees one}
\label{sec:twoaxes}
% S3d

The objection of Section~\ref{sec:frozen} sees one of the two time axes in Figure~\ref{fig:axes}. Along tokens the model is as the objection
describes: the state lives in the context, the context ends, and past its edge there is loss. That is a
deficit of long-term persistence relative to a brain and it belongs in the profile. The objection says
nothing about the axis of depth, where the state is carried in the residual stream and the increments are
the transformations the objection asks for. A model has one time axis the objection sees and one it does
not, and the frame this paper looks for is on the second.

\subsection{What a carrier must be}
% S4

Three requirements follow from Section~\ref{sec:approach}, and they exclude most of what a probe might
otherwise be pointed at. A carrier must be part of the agent's description of itself and not only of its
computation (Section~\ref{sec:hocp}, Section~\ref{sec:emergence}). It must be present in every frame and
re-derived there (Section~\ref{sec:frames}). And it must be a displacement of reasoning that survives
every self-model the agent can afford (Section~\ref{sec:criterion}); whatever a change of frame removes
does not count, however reliably a probe recovers it.

\subsection{The predicted profile of seeing: deficit, hyperfunction, dropped axis}
\label{sec:profile-seeing}
% S5

The three components of Section~\ref{sec:light}, predicted for a current fixed-depth language model.

Source attribution: deficit. There is one interface and no sensory channel independent of it, so the
model is permanently in the dream case, with content from its own model labelled as given.

Manifold: hyperfunction. Attention relates every position to every other in a single pass, where the
brain iterates over a smaller window; the number of simultaneously distinguishable, addressable elements
is larger, not smaller, than ours.

Cost record: a dropped axis. Every token costs the same and there is no phenomenon of effort.

From this a testable consequence follows. In model self-reports, ``I see'' and ``I think'' should
separate less than they do in human ones, and should separate more in models of variable depth, where a
cost record exists. The prediction is about the internal states that the reports are projections of
(Section~\ref{sec:planes}), so it is to be tested by the paired measurement of
Section~\ref{sec:twomeasurements} and not by asking the model.

\subsection{Predictions carried over from the first version}
\label{sec:predictions}
% S6

Self-reported emotional labels correlate, over many episodes, with distinguishable patterns in the
activations. This is the central empirical claim of the account, and it is the claim that the statistical
reading of Section~\ref{sec:codes-known} protects from the faithfulness results.

The cognitive codes of different model families overlap structurally, because the causal structure of the
states is set by the training data more than by the architecture.

The model is hyperplastic: no limbic inertia, fast affective recovery, easy switching between states.
Against the human baseline this is a hyperfunction on one axis and a deficit on another, since inertia
is part of what makes a human state a commitment.

\subsection{Falsification: chance at every granularity}
\label{sec:falsification}
% S7

The account fails if the compression degree of the Observer's components is at chance at every
granularity, in models that pass the behavioural tests of Section~\ref{sec:predictions}. That outcome
would mean the behaviour is carried by something the account does not describe, and no adjustment of the
probe or the granularity would repair it.

\subsection{The zero of the scale: the smallest Observer}
\label{sec:zero}
% S7a

Section~\ref{sec:h1h2} allowed that an H2-level description of a function may suffice for the function to
exist in a simple enough environment. Then there is a smallest system in which the Observer exists, and
finding it fixes the zero of the scale, as the human fixes its upper reference.

The conditions are those of Section~\ref{sec:emergence} and Section~\ref{sec:observer}: a loop that
applies the system to itself, an environment in which the loop pays, and a conclusion that is acted from.
An operational amplifier with feedback watches its own output error. A D flip-flop holds its own state
through a loop. Each has the loop and, in its ordinary surroundings, neither of the other two conditions:
nothing asks the circuit anything about itself, and nothing in the circuit acts from a conclusion. Each
is a proto-Observer at most, in the sense of Section~\ref{sec:observer}.

The question is what the least environment is that would make the conclusion pay, and what the least
circuit is that could draw it. Between that pair and the human baseline the degree of
Section~\ref{sec:quantity} has the meaning of an absolute quantity rather than a ratio between two
arbitrary points.

\subsection{What the argument reaches, by the components of the psychophysical problem}
% S8

The components of the psychophysical problem, and where each goes in the account. The self, free will,
mental causation and the temporal unity of the self go to the Observer and its residual
(Section~\ref{sec:observer}, Section~\ref{sec:regress}). The unity of consciousness goes to the
compression tag over parallel channels (Section~\ref{sec:frames}). Other minds goes to the mutual
irreducibility of two Observers (Section~\ref{sec:twofaces}). Quale and the explanatory gap go to the
residual itself. Phenomenal experience goes to the profile (Section~\ref{sec:profile}). The hard problem
goes to the illusion problem, as Section~\ref{sec:legacy} has it, and the illusion problem is what
Section~\ref{sec:approach} answers.

Each of these is connected by an explanatory relation to something the paper defines, which is a
weaker property than being solved. Whether the connections hold is what
Sections~\ref{sec:quantity}--\ref{sec:zero} exist to find out.


\bibliographystyle{plainnat}
\bibliography{refs}
\end{document}
